%%
%% fall-lecture12.tex
%% 
%% Made by Alex Nelson <pqnelson@gmail.com>
%% Login   <alex@lisp>
%% 
%% Started on  2026-10-24T12:20:40-0700
%% Last update 2026-10-24T12:20:40-0700
%% 

\lecture[Cohomology Ring of Infinite Grassmann Manifold]{}

\begin{recall}
We have $G_{n}=G_{n}(\RR^{\infty})$ consist of all $n$-dimensional
subspaces of $\RR^{\infty}$. We endowed it with a direct limit
topology.

We also saw $\gamma^{^n}$ the $\RR^{n}$-bundle over $G_{n}$ defined as
\begin{equation}
\gamma^{n}=\{(X,v)\in G_{n}\times\RR^{\infty}\mid v\in X\}.
\end{equation}
For now, let us assume that $w_{1}(\gamma^{n})$, \dots,
$w_{n}(\gamma^{n})$ exists.
\end{recall}

\begin{theorem}\label{thm:math157A:lecture12:cohomology-of-grass}
The cohomology ring $H^{*}(G_{n};\ZZ_{2})$ is a polynomial ring over
$\ZZ_{2}$ with generators $w_{1}(\gamma^{n})$, \dots,
$w_{n}(\gamma^{n})$; i.e.,
\begin{equation}
H^{*}(G_{n};\ZZ_{2})\iso\ZZ_{2}[w_{1}(\gamma^{n}),\dots,w_{n}(\gamma^{n})].
\end{equation}
\end{theorem}

To prove this, we will show $H^{*}(G_{n};\ZZ_{2})$ contains a subring
isomorphic to
$\ZZ_{2}[w_{1}(\gamma^{n}),\dots,w_{n}(\gamma^{n})]$. Then we will
show there is nothing else in the cohomology ring of $G_{n}$.

\begin{lemma}
There are no non-zero polynomial relations between $w_{1}(\gamma^{n})$, \dots,
$w_{n}(\gamma^{n})$. In other words, there does not exist a nonzero
$p(x_{1},\dots,x_{n})\in\ZZ_{2}[x_{1},\dots,x_{n}]$ such that
$p(w_{1}(\gamma^{n}),\dots,w_{n}(\gamma^{n}))=0$.
\end{lemma}

\begin{proof}
Consider $\xi$ an $\RR^{n}$-bundle over a paracompact space. Then we
can always find some map $f\colon B(\xi)\to G_{n}$ such that
the pullback bundle $f^{*}\gamma^{n}\iso\xi$ are isomorphic as
bundles. If we had some polynomials relationship, then by the axioms:
\begin{subequations}
  \begin{align}
p(w_{1}(\gamma^{n}),\dots,w_{n}(\gamma^{n}))
&= p(f^{*}w_{1}(\xi),\dots,f^{*}w_{n}(\xi))\\
&=f^{*}p(w_{1}(\xi),\dots,w_{n}(\xi)).
  \end{align}
\end{subequations}
So it suffices to prove there is no polynomial relationship. Our goal
is now to construct such an $\RR^{n}$-bundle $\xi$.

Let us try $\xi=\gamma^{1}\oplus\dots\oplus\gamma^{1}$ the Whitney sum
of $n$ copies of $\gamma^{1}$. This is certainly an $\RR^{n}$-bundle
over $\RP^{\infty}$ (which is good), and
\begin{equation}
w(\gamma^{1}\oplus\dots\oplus\gamma^{1})=(1+a)^{n}.
\end{equation}
Now, observe $w_{0}=1$, $w_{1}=na$, $w_{2}=\frac{1}{2}n(n-1)a^{2}$,
\dots, $w_{n}=a^{n}$. Already, we see there exists a polynomial
relationship between $w_{1}$ and $w_{2}$ since: for even $n$ we have
$w_{1}=0$, and for odd $n$ we have $w_{2}=w_{1}^{2}$. So this is not
quite enough to use the Whitney sum of \emph{the same} $\gamma^{1}$.

The trick is to consider the projection maps
\begin{equation}
\pi_{j}\colon\underbrace{\RP^{\infty}\times\cdots\times\RP^{\infty}}_{n~\text{copies}}\to\RP^{\infty}
\end{equation}
for $j=1,\dots,n$. Then we define
\begin{equation}
\xi = \pi_{1}^{*}\gamma^{1}\oplus\dots\oplus\pi_{n}^{*}\gamma^{1}
\end{equation}
as the Whitney sum of the \emph{pullbacks} of $\gamma^{1}$ under the
projection maps. This again gives us an $\RR^{n}$-bundle over
$(\RP^{\infty})^{n}$.

Let us determine the Stiefel--Whitney classes. We see
\begin{equation}
w(\xi)=\left(\pi_{1}^{*}w(\gamma^{1})\right)(\dots)\left(\pi_{n}^{*}w(\gamma^{1})\right)
\end{equation}
where each $w(\gamma^{1})=1+a$. But the pullback morphism induces a
morphism of the cohomology rings
\begin{equation}
\pi_{j}^{*}\colon H^{*}(\RP^{\infty};\ZZ_{2})\to H^{*}((\RP^{\infty})^{n};\ZZ_{2})\iso
\underbrace{H^{*}(\RP^{\infty};\ZZ_{2})\otimes\dots\otimes H^{*}(\RP^{\infty};\ZZ_{2})}_{n~\text{times tensor product}}
\end{equation}
by Kunneth formula (since $\ZZ_{2}$ is torsion-free), which sends
\begin{equation}
a\mapsto(\underbrace{1\otimes\dots\otimes1}_{j-1}\otimes a\otimes1\otimes\dots\otimes1)=:a_{j}.
\end{equation}
Then when we pullback, we get $\pi_{j}^{*}(w(\gamma^{1}))=1+a_{j}$, so 
\begin{equation}
w(\xi)=(1+a_{1})(1+a_{2})(\dots)(1+a_{n}),
\end{equation}
and each $a_{j}$ lives in different copies of the cohomology. We see
that
\begin{subequations}
\begin{equation}
w_{1}=a_{1}+\dots+a_{n},
\end{equation}
\begin{equation}
w_{2}=\sum_{i<j}a_{i}a_{j},
\end{equation}
and so on until we get
\begin{equation}
w_{n}=a_{1}a_{2}(\cdots)a_{n}.
\end{equation}
\end{subequations}
Observe $w_{j}$ is just the $j^{\text{th}}$ elementary symmetric
polynomial with variables $a_{1}$, $a_{2}$, \dots, $a_{n}$.

\textsc{Algebraic fact}: \emph{there are no polynomial relations
between elementary symmetric polynomials}.

This means there are no polynomial relations between the
$w_{i}(\xi)$. And it concludes the proof of the Lemma.
\end{proof}

\begin{proof}[{Proof (Theorem~\ref{thm:math157A:lecture12:cohomology-of-grass})}]
Let us count the dimension $\dim_{\ZZ_{2}}H^{r}(G_{n};\ZZ_{2})$. We
see that
$|\{(b_{1},\dots,b_{n})\in\NN_{0}^{n}\mid\sum_{k}kb_{k}=n\}|$ is a
lower bound on the dimension because it coincides with degree-$r$
monomials in unknowns $w_{1}(\gamma^{n})$, \dots,
$w_{n}(\gamma^{n})$. An upper bound is
\begin{equation}
\dim_{\ZZ_{2}}H^{r}(G_{n};\ZZ_{2})\leq\#\mbox{($r$-cells in CW complex of $G_{n}$)}.
\end{equation}
Then, in Chapter 6 of the textbook~\cite{milnor1974characteristic},
they prove the number of $r$-cells in the CW structure of $G_{n}$ is
precisely 
$|\{(b_{1},\dots,b_{n})\in\NN_{0}^{n}\mid\sum_{k}kb_{k}=n\}|$.
Hence the theorem.
\end{proof}

\begin{remark}
Reflect on the proof of the Lemma.
Do we \emph{need} the universal property of $\gamma^{n}$?
Although it \emph{appears} necessary, we do not need it. We have
\begin{equation}
f\colon(\RP^{\infty})^{n}\to G_{n}
\end{equation}
such that $f^{*}\gamma^{n}\iso\xi=\pi_{1}^{*}\gamma^{1}\oplus\dots\oplus\pi_{n}^{*}\gamma^{1}$.
If we could explicitly construct such an $f$, then we do not need to
rely on the universal property.

We can map $(\ell_{1},\dots,\ell_{n})\mapsto J(\ell_{1}\oplus\dots\oplus\ell_{n})$
where each $\ell_{j}$ is a line in $\RR^{\infty}$, but
$\ell_{1}\oplus\dots\oplus\ell_{n}$ is a line in
$(\RR^{\infty})^{n}$. So we use the isomorphism
\begin{equation}
J\colon(\RR^{\infty})^{n}\to\RR^{\infty}.
\end{equation}
We still need to show $f^{*}\gamma^{n}\iso\xi$.

Let us fix a point $(\ell_{1},\dots,\ell_{n})\in\RP^{\infty}$, the
fiber of $f^{*}\gamma^{n}$ (is just the fiber of $\gamma^{n}$) is
$J(\ell_{1}\oplus\dots\oplus\ell_{n})$; and the fiber in $\xi$ is just
$\ell_{1}\oplus\dots\oplus\ell_{n}$ the direct sum. So $J$ gives the
isomorphism.

Once we have our $f$, we can look at the pullback of the cohomology
$f^{*}\colon H^{*}(G_{n};\ZZ_{2})\to
H^{*}((\RP^{\infty})^{n};\ZZ_{2})$ which is injective. This follows
from the preceding construction.
\end{remark}

\begin{theorem}[Uniqueness of Stiefel--Whitney classes]
The assignment $\xi\mapsto w(\xi)$ satisfying all four Axioms for
Stiefel--Whitney classes is unique.
\end{theorem}

\begin{proof}
Suppose we have two such assignments $w(\xi)$ and $\widetilde{w}(\xi)$ 
satisfying all four axioms. The last axiom tells us
\begin{equation}
w(\gamma^{1}_{1})=1+a=\widetilde{w}(\gamma^{1}_{1}).
\end{equation}
So for $\gamma^{1}_{1}$, we have the following commutative diagram
\begin{equation}
\vcenter{\xymatrix{\gamma_{1}^{1}\ar[r]\ar[d] & \gamma^{1}\ar[d]\\
\RP^{1}\ar[r] & \RP^{\infty}}}
\end{equation}
which implies for the universal line bundle
\begin{equation}
w(\gamma^{1})=\widetilde{w}(\gamma^{1})=1+a.
\end{equation}
We can choose $\xi=\pi_{1}^{*}\gamma^{1}\oplus\dots\oplus\pi_{n}^{*}\gamma^{1}$
and by the preceding remark we have a map
\begin{equation}
f\colon(\RP^{\infty})^{n}\to G_{n}
\end{equation}
such that $f^{*}\gamma^{n}\iso\xi$ and injectivity on cohomology ring
by the previous theorem. Then
\begin{equation}
w(\xi)=(1+a_{1})(\dots)(1+a_{n})=\widetilde{w}(\xi).
\end{equation}
Then if we pullback $f^{*}w(\gamma^{n})=w(\xi)$ and
$f^{*}\widetilde{w}(\gamma^{n})=\widetilde{w}(\xi)$, hence
$w(\gamma^{n})=\widetilde{w}(\gamma^{n})$. 
\end{proof}